% =============================================================================
\section{Introduction}\label{sec:intro}
% =============================================================================

\subsection{Causality before geometry}

General relativity treats spacetime geometry as dynamical but primary; quantum field theory treats it
as fixed and equally primary. In both, matter and fields evolve on a manifold that exists before
anything happens on it. Causal Physics adopts the opposite stance: the primitive of physics is a
set of events and the relation of influence between them,
\begin{equation}
  u \prec v \quad\Longleftrightarrow\quad u \text{ can influence } v ,
\end{equation}
and geometry is a large-scale description of regularities in that relation.

This stance is not new. It is the founding hypothesis of causal set theory, compressed by Sorkin
into the slogan ``order $+$ number $=$ geometry''~\cite{Bombelli1987,Surya2019}: a causal order
determines a Lorentzian metric up to a conformal factor, and the number of events in a region
fixes that factor. What Causal Physics adds is a physical reading of the order. Events are updates
of a substrate; persistent flows of events are matter; and gravity is the way the order itself
responds to those flows, slowing clocks, narrowing causal cones and adding events near matter. The
question this paper asks is how much of gravitation follows from that picture once it is made
precise enough to be computed.

\subsection{Method: a laboratory that tests and does not rescue}\label{sec:method}

All quantitative statements in this paper come from the Causality Lab~\cite{CausalityLab}, an open
collection of numerical tests written for this programme. The rules of the laboratory matter as
much as its results, because a theory developed together with its own tests can easily test only
what it already put in.

Every test states its question, prediction, null model and numerical criterion before the full run,
declares the level of conclusion it can support, and ends with one of four verdicts:
\verdict{passed}, \verdict{failed}, \verdict{undecided} or \verdict{instrument failed}. Negative
results are kept, and a failed test is never rescued by changing its criteria. The rules are stated
in full as the methodological principles M1--M7 (Sec.~\ref{sec:method-ax}).

Throughout the paper, results are labelled with the identifiers of the tests (e.g.\ F7.4a) so that
each can be rerun. Appendix~\ref{app:lab} lists the tests used.

\subsection{Main results}

\begin{enumerate}
\item \textbf{The substrate} is a Lorentz-invariant causal set; a finite-degree substrate with a rest
  frame is excluded by the tests together with a known theorem (Sec.~\ref{sec:substrate}).
\item \textbf{Geometry is read blindly from the order.} An instrument that receives only events
  and their relation recovers the clock potential $\phi$ and the spatial potential $\psi$ of a weak
  field, and with them the post-Newtonian parameter $\gamma$, in 2+1 and 3+1 dimensions
  (Sec.~\ref{sec:reading}).
\item \textbf{Matter carries momentum.} A flow that keeps its timelike direction as a state stays
  localised and moves with its source, and its energy--momentum tensor is measured correctly
  (Sec.~\ref{sec:flows}).
\item \textbf{The response of the order fixes $\gamma$.} A fixed order gives $\gamma=-1$, narrowing
  cones alone give $\gamma=1/d_s$, and conservation of causal throughput gives
  $\gamma=2/(d_s-1)$; on 3+1 sprinklings the last reads $\gamma=1.000\pm0.045$ (raw value; the
  formal verdict is undecided because the comparison metric was read with a larger error,
  Sec.~\ref{sec:response}).
\item \textbf{The coupling is tensorial.} In 3+1 the throughput principle is the linearised coupling
  of general relativity to $T_{ab}$; a substrate rule realising it makes pressure and motion
  gravitate and drags the cones (Part~\ref{part:geometric}).
\item \textbf{The rule follows from the action of the order.} The expected Benincasa--Dowker action
  of a causal set carries the Einstein--Hilbert stiffness in the static conformal sector,
  $6.0003$ against $6$, in the mixed sector, $7.983$ against $8$, and in the optical sector, $2.002$
  against $2$; $\gamma=1$ follows ($0.997$ to $1.000$). The discreteness adds a non-analytic correction and
  non-local age terms (Part~\ref{part:geometric}).
\item \textbf{The road to predictions} is the static quadratic form of the order, now known in all three
  sectors; the asymptotic age term of the optical sector is the next computation
  (Part~\ref{part:predictions}).
\item \textbf{Where to look.} The age term is not suppressed like the local corrections. For a cosmological
  age and $\ell_k$ at the level of existing bounds it places the breakdown of the static response at tens of
  micrometres, the range of current torsion balances (Sec.~\ref{sec:window}).
\end{enumerate}

\subsection{Relation to existing frameworks}

\paragraph{Causal set theory.} The substrate, the reading of geometry from order and number and the
Benincasa--Dowker operator and action are taken from causal set
theory~\cite{Bombelli1987,Surya2019,BenincasaDowker2010,Sorkin2007,DowkerGlaser2013}. The
contribution of this work lies elsewhere: in the physical reading of matter as flows that carry
momentum, in the principle by which the order responds to them, in the explicit second-order
computation of the expected action with its non-local age term, and in the systematic numerical
testing of each step.

\paragraph{General relativity.} The theory does not compete with general relativity in the regime
where general relativity is tested. In the linear regime it reproduces it, and beyond the linear
order it adopts it, with the intention of deriving it thermodynamically along the lines of
Jacobson~\cite{Jacobson1995}. Its own content lies in the discreteness and non-locality of the
substrate.

\paragraph{Uniqueness theorems.} A massless spin-2 field coupled consistently to the total conserved
energy--momentum is general relativity~\cite{Deser1970,FeynmanGravitation}, and in four dimensions
the Einstein tensor is the unique divergence-free tensor built from the metric and its first two
derivatives~\cite{Lovelock1971}. The results of Part~\ref{part:geometric} are consistent with these
theorems and use them: once the order is shown to couple to $T_{ab}$ with the Einstein--Hilbert
stiffness, the non-linear completion is fixed.

\paragraph{Modified gravity.} The theory makes no claim of a modified low-acceleration regime
(Sec.~\ref{sec:astro-cosmo}).

\subsection{Structure of the paper}

Part~\ref{part:ontology} describes, without equations, how Causal Physics sees the world: events,
order and time, observation, transport, merging, matter, gravity, horizons and superposition, each
statement labelled by how firmly it stands. Part~\ref{part:foundations} sets up the substrate, the reading of geometry from it, matter as flows
and the response of the order. Part~\ref{part:units} states the parameters and the conditional
bridge to SI units. Part~\ref{part:axioms} collects the axioms in their revised form, with the status
of each. Part~\ref{part:geometric} derives and tests the gravitational coupling, from the throughput
principle to the Einstein--Hilbert stiffness of the discrete action. Part~\ref{part:astro} states what
remains of the strong-field and astrophysical claims. Part~\ref{part:predictions} sets out the
predictions the theory can make and how they can fail. Part~\ref{part:discussion} discusses the
limitations. Appendix~\ref{app:lab} describes the laboratory and the tests;
Appendix~\ref{app:changes} relates this formulation to the first published version of the
programme.
